\documentclass[10pt,a4paper]{amsart}
\usepackage[margin=19mm]{geometry}
\usepackage{amsmath,amssymb,mathtools}
\usepackage[scr=rsfs,cal=euler]{mathalfa}
\usepackage{xfrac}
\usepackage[unicode,hidelinks,bookmarks=false]{hyperref}
\newcommand{\ie}{i.e.\ }
\newcommand{\cf}{cf.\ }
\newcommand{\resp}{respectively\ }
\newcommand{\Subsection}{\subsection}
\providecommand{\nobreakdash}{\nobreak\hbox{-}}
\setlength{\emergencystretch}{2em}
\multlinegap0pt
\usepackage{amssymb}
\usepackage{mathtools}
\usepackage[shortlabels]{enumitem}
\newtheorem{definition}{Definition}
\newtheorem{lemma}{Lemma}
\newcommand{\En}{\mathcal E_n}
\newcommand{\esort}{\operatorname{esort}}
\title{Quadratic Gr\"obner bases for cut ideals of cycles and ring graphs}
\author{Hidefumi Ohsugi}
\date{Source: arXiv:2609.11512v1 (CC BY 4.0)}
\begin{document}
\maketitle

\noindent\emph{Source and licence.} Quadratic Gr\"obner bases for cut ideals of cycles and ring graphs. Version arXiv:2609.11512v1 (CC BY 4.0), distributed by arXiv under the Creative Commons Attribution 4.0 International licence, http://creativecommons.org/licenses/by/4.0/, as recorded in the arXiv licence metadata for this exact version and shown on the abstract page https://arxiv.org/abs/2609.11512v1. Full licence notice: \texttt{licenses/arxiv-2609.11512-CC-BY-4.0.txt}.\par\smallskip

\noindent\textit{Editorial scope.} Counted unit: the article's lemma lem:quadraticminimum (source0.tex lines 465--467). The article's complete original proof of it is delivered with it (source0.tex lines 469--866, one \texttt{proof} environment of 398 lines). No mathematics is written, repaired or re-derived: the statement and the proof are the article's own lines, in the article's own notation, and no retained line is substituted or re-flowed.  Retained uncounted as the source-local context the proof needs: lines 280--307.  Nothing was omitted from any retained span, so the package carries no omission record.  No retained line contains a citation, so the wrapper carries no bibliography and no bibliography entry.  No retained line contains a figure environment, an image-inclusion command or drawing code, so no image is delivered and the package declares no asset dependency.  Editorial formatting only, in addition: an AMS \texttt{amsart} preamble that restates the article's own declarations and macros used by the retained lines, namely the environments \texttt{definition}, \texttt{lemma} on separate counters, together with the standard packages those lines need.  The retained environments print in this document as Definition 1, Lemma 1 in order of appearance, whereas the article prints the same items with the numbering of its own full document.  The complete native source of the article as distributed in the arXiv e-print for this exact version is bundled unchanged as \texttt{sources/210026/source0.tex}.
\begin{definition} \label{def:evensort}
Fix $A,B\in\En$. 
Let
\[
I=A\cap B,
 \qquad
A\triangle B=\{d_1<d_2<\cdots<d_{2m}\},
 \qquad
p\in\{0,1\} \mbox{ with } p\equiv |I|\pmod2.
\]
Since $A$ and $B$ are both even,
the symmetric difference $A \triangle B$ has even (possibly zero) cardinality.
\begin{enumerate}
    \item 
If $m\equiv p\pmod2$, define
\begin{align*}
 A^\sharp&=I\cup\{d_{2r-1}:1\le r\le m\},\\
 B^\sharp&=I\cup\{d_{2r}:1\le r\le m\}.
\end{align*}
\item
If $m\not\equiv p\pmod2$ (which necessarily implies $m\ge1$), define
\begin{align*}
 A^\sharp&=I\cup\{d_{2r-1}:1\le r\le m\}\cup\{d_{2m}\},\\
 B^\sharp&=I\cup\{d_{2r}:1\le r\le m-1\}.
\end{align*}
\end{enumerate}
Here a set indexed by an empty range is understood to be empty. We call the unordered pair $\esort(A,B)=\{A^\sharp,B^\sharp\}$ the \emph{even-sorted pair} associated with $A,B$. We say that $\{A,B\}$ is even-sorted if $\{A,B\}=\esort(A,B)$.
\end{definition}

\begin{lemma}\label{lem:quadraticminimum}
In every quadratic fiber, the even-sorted monomial uniquely minimizes $\Omega(U)+\Omega(V)$.
\end{lemma}

\begin{proof}
Fix a quadratic fiber
\[
 \mathcal{F}^{(2)}_s
=
 \left\{
q_Uq_V : U,V\in\mathcal{E}_n,\;
 \chi_U+\chi_V=s
 \right\},
 \qquad s\in\{0,1,2\}^n.
\]
Set $I_s:=\{i\in[n]:s_i=2\}$ and $D_s:=\{i\in[n]:s_i=1\}$. For every $q_Uq_V\in\mathcal{F}^{(2)}_s$, one has $U\cap V=I_s$ and $U\triangle V=D_s$. 
Let $D_s=\{d_1<d_2<\cdots<d_{2m}\}$.
If $D_s=\emptyset$, then $U=V=I_s$ for every $q_Uq_V\in\mathcal F_s^{(2)}$. Hence the fiber consists of the single monomial $q_{I_s}^2$, and the assertion is immediate. We may therefore assume that $D_s\ne\emptyset$, or equivalently, that $m\ge 1$.
Let $p\in\{0,1\}$ be determined by $p\equiv |I_s|\pmod 2$, and let $q_{U^\sharp}q_{V^\sharp}$ be the even-sorted monomial determined by $I_s$ and $D_s$ as in Definition~\ref{def:evensort}. 
This monomial depends only on the fiber.


For $1\le k\le n$, let
\begin{align}   \label{eq:s and delta}
S_k:=F_k(U)+F_k(V),
\qquad
\Delta_k:=F_k(U)-F_k(V).
\end{align}
Since $\chi_U+\chi_V=s$, we have
\[
S_k
=
F_k(U)+F_k(V)
=
\sum_{i=1}^k\bigl(\chi_U(i)+\chi_V(i)\bigr)
=
\sum_{i=1}^k s_i.
\]
Thus $S_k$ depends only on the fiber
$\mathcal F_s^{(2)}$.
From \eqref{eq:s and delta},
it follows that
\[
F_k(U)^2+F_k(V)^2
=
\frac12\bigl(S_k^2+\Delta_k^2\bigr).
\]
Hence
\[
P(U)+P(V)
=
\frac12\sum_{k=1}^n
\bigl(S_k^2+\Delta_k^2\bigr).
\]
Similarly,
\[
L(U)+L(V)
=
\sum_{k=1}^n S_k,
\qquad
L(U)-L(V)
=
\sum_{k=1}^n \Delta_k,
\]
and hence
\[
L(U)^2+L(V)^2
=
\frac12\left(
\left(\sum_{k=1}^n S_k\right)^2
+
\left(\sum_{k=1}^n\Delta_k\right)^2
\right).
\]
Consequently,
\[
 \Omega(U)+\Omega(V)
= C_s+\frac{M}{2}Q(U,V)-\frac12R(U,V)^2,
\]
where
\[
 C_s:=\frac M2\sum_{k=1}^nS_k^2
 -\frac12\left(\sum_{k=1}^nS_k\right)^2,
\qquad
Q(U,V):=\sum_{k=1}^n\Delta_k^2,
 \qquad
R(U,V):=\sum_{k=1}^n\Delta_k.
\]
Note that $C_s$ depends only on the fiber.
Since $|\Delta_k|\le k$, we have
\[
 |R(U,V)|\le\sum_{k=1}^n|\Delta_k|
 \le\sum_{k=1}^n k=N.
\]

Let $ q_Uq_V,\ q_{U'}q_{V'}\in\mathcal{F}^{(2)}_s$.
\begin{itemize}
    \item 
If $Q(U,V)<Q(U',V')$, then 
$Q(U',V') - Q(U,V) \ge 1$
and hence
we have
\begin{align*}
& \quad ( \Omega(U')+\Omega(V'))-( \Omega(U)+\Omega(V))\\
& =
 \frac{M}{2}\bigl(Q(U',V')-Q(U,V)\bigr)
-
 \frac12\bigl(R(U',V')^2-R(U,V)^2\bigr)\\
& \ge
 \frac{M}{2}-\frac{N^2}{2}>0
\end{align*}
since $|R(U,V)|, |R(U',V')| \le N$.
Thus every minimizer of \(\Omega(U)+\Omega(V)\) must first minimize
\(Q(U,V)\).
\item
If $Q(U,V)=Q(U',V')$, then 
\[
( \Omega(U')+\Omega(V'))-( \Omega(U)+\Omega(V))
=-
 \frac12\bigl(R(U',V')^2-R(U,V)^2\bigr).
\]
Thus among the minimizers of \(Q\), minimizing
\(\Omega(U)+\Omega(V)\) is equivalent to maximizing
$|R(U,V)|$.
\end{itemize}

We now determine the minimizers of \(Q\). 
Recall that $D_s=U\triangle V=\{d_1<\cdots<d_{2m}\}$.
Note that $\Delta_k$ can change only when $k\in D_s$.
Hence $\Delta_k$ is constant between two consecutive elements of $D_s$.
Let $\{U^\sharp,V^\sharp\}=\esort(U,V)$ be the even-sorted pair determined by the fiber. 
We may assume that \(d_1\in U^\sharp\).
Let $ \Delta_k^\sharp
:= F_k(U^\sharp)-F_k(V^\sharp)$.

\medskip

\noindent
{\bf Case 1} (\(m\equiv p\pmod2\)). Then $|U^\sharp|=|V^\sharp|$
and $\Delta_k^\sharp$
is given explicitly by
\[
 \Delta_k^\sharp=
 \begin{cases}
1, & d_{2r-1}\le k<d_{2r}
    \text{ for some }r,\\
0, & \text{otherwise}.
 \end{cases}
\]
In particular, $(\Delta_k^\sharp)^2=\Delta_k^\sharp $.
Let $\ell_r:=d_{2r}-d_{2r-1}>0$. 
It then follows that
\[
Q(U^\sharp,V^\sharp) =
R(U^\sharp,V^\sharp) =
 \sum_{r=1}^m\ell_r.
\]

Let $q_Uq_V\in\mathcal{F}^{(2)}_s$.
If $d_{2r-1}\le k<d_{2r}$, then exactly $2r-1$ elements of
$D_s=U\triangle V$ belong to $[k]$. 
Hence
\[
\Delta_k
=
|(U\setminus V)\cap[k]|
-
|(V\setminus U)\cap[k]|
\equiv
|(U\setminus V)\cap[k]|
+
|(V\setminus U)\cap[k]|
= 2r - 1
\pmod 2
,\]
and therefore $|\Delta_k|\ge1$.
Consequently,
\[
Q(U,V) =
 \sum_{k=1}^n\Delta_k^2
 \ge
 \sum_{r=1}^m
 \sum_{k=d_{2r-1}}^{d_{2r}-1}1
=
 \sum_{r=1}^m\ell_r
= Q(U^\sharp,V^\sharp).
\]
Hence the even-sorted pair minimizes $Q$.

Moreover, equality
\[
Q(U,V)=Q(U^\sharp,V^\sharp)
\]
can hold if and only if 
\[
 |\Delta_k|=
 \begin{cases}
1, & d_{2r-1}\le k<d_{2r}
    \text{ for some }r,\\
0, & \text{otherwise}.
 \end{cases}
\]
In such a case,
there are signs $\sigma_r\in\{\pm1\}$
such that
\[
 \Delta_k=\sigma_r
 \qquad
(d_{2r-1}\le k<d_{2r}).
\]
Hence every \(Q\)-minimizer satisfies
\[
R(U,V) =
 \sum_{r=1}^m\sigma_r\ell_r.
\]
By the triangle inequality,
\[
|R(U,V)|
 \le
 \sum_{r=1}^m\ell_r
= |R(U^\sharp,V^\sharp)|.
\]
Since every \(\ell_r\) is positive, equality holds if and only if
\[
 \sigma_1=\sigma_2=\cdots=\sigma_m.
\]
The two possible common signs correspond precisely to interchanging
\(U^\sharp\) and \(V^\sharp\).
Among the $Q$-minimizers, $|R|$ is therefore maximal only for $\{U^\sharp,V^\sharp\}$. Hence $q_{U^\sharp}q_{V^\sharp}$ is the unique quadratic minimizer in Case~1.

\medskip

\noindent
{\bf Case 2} (\(m\not\equiv p\pmod2\)). 
Then
\begin{align*}
U^\sharp
&=I_s\cup\{d_{2r-1}:1\le r\le m\}\cup\{d_{2m}\},\\
V^\sharp
&=I_s\cup\{d_{2r}:1\le r\le m-1\}.
\end{align*}
Hence $\Delta_k^\sharp$
is given by
\[
 \Delta_k^\sharp=
 \begin{cases}
1, & d_{2r-1}\le k<d_{2r}
    \text{ for some }r,\\
0, & d_{2r}\le k<d_{2r+1}
    \text{ for some }r<m,\\
2, & d_{2m}\le k\le n.
 \end{cases}
\]
Let $\ell_r:=d_{2r}-d_{2r-1}>0$ and $\tau:=n-d_{2m}+1>0$. It follows that
\[
Q(U^\sharp,V^\sharp) =
 \sum_{r=1}^m\ell_r+4\tau,
 \qquad
R(U^\sharp,V^\sharp) =
 \sum_{r=1}^m\ell_r+2\tau.
\]

Let $q_Uq_V\in\mathcal{F}^{(2)}_s$.
From the same discussion in Case 1,
on every interval
$d_{2r-1}\le k<d_{2r}$,
we have $|\Delta_k|\ge1$.
There is an additional restriction after the last element
\(d_{2m}\).  Since both \(|U|\) and \(|V|\) are even,
\[
|U\setminus I_s|\equiv |V\setminus I_s|
 \equiv |I_s|
 \equiv p
 \pmod2.
\]
On the other hand,
\[
|U\setminus I_s|+|V\setminus I_s|=2m.
\]
Moreover,
\[
 \Delta_n
=
F_n(U) - F_n(V)
= |U|-|V| 
= |U \setminus I_s|-|V \setminus I_s| 
\equiv |U \setminus I_s|+|V \setminus I_s| 
=2m
\equiv0
\pmod 2.
\]
If \(\Delta_n=0\), then
\[
|U\setminus I_s|=m,
\]
which would imply
\[
m\equiv p\pmod2,
\]
contrary to the assumption of Case~2. 
Since $\Delta_n$ is even,
we have $|\Delta_n|\ge2$.
Since no element of the symmetric difference occurs after \(d_{2m}\), the value of \(\Delta_k\) is constant for \(k\ge d_{2m}\). Thus
\[
|\Delta_k|\ge2
 \qquad(d_{2m}\le k\le n).
\]
Consequently,
\[
Q(U,V) =
 \sum_{k=1}^n\Delta_k^2
\ge
 \sum_{r=1}^m
 \sum_{k=d_{2r-1}}^{d_{2r}-1}1
+
 \sum_{k=d_{2m}}^n4
=
 \sum_{r=1}^m\ell_r+4\tau
= Q(U^\sharp,V^\sharp).
\]
Hence the even-sorted pair minimizes $Q$.

Suppose now that equality 
\[
Q(U,V)=Q(U^\sharp,V^\sharp)
\]
holds.
Then 
\[
 |\Delta_k|=
 \begin{cases}
1, & d_{2r-1}\le k<d_{2r}
    \text{ for some }r,\\
0, & d_{2r}\le k<d_{2r+1}
    \text{ for some }r<m,\\
2, & d_{2m}\le k\le n.
 \end{cases}
\]
In such a case,
there are signs $\sigma_r\in\{\pm1\}$
such that
\[
 \Delta_k=\sigma_r
 \qquad
(d_{2r-1}\le k<d_{2r}).
\]
In particular,
$$
 \Delta_{d_{2m}-1}=\sigma_m.
$$
Since exactly one of \(U\) and \(V\) contains \(d_{2m}\), 
we have
$$
 \Delta_{d_{2m}}-\Delta_{d_{2m}-1}\in\{1,-1\}.
$$
On the other hand, equality in the lower bound for \(Q\) requires
$$
 |\Delta_{d_{2m}}|=2.
$$
Since \(\Delta_{d_{2m}-1}=\sigma_m\in\{1,-1\}\), the only possibility is
$$
 \Delta_{d_{2m}}=2\sigma_m.
$$
Finally, since \(d_{2m}\) is the largest element of \(U\triangle V\), \(\Delta_k\) does not change for \(k\ge d_{2m}\). Hence
$$
 \Delta_k=2\sigma_m
 \qquad
 (d_{2m}\le k\le n).
$$
Therefore every \(Q\)-minimizer satisfies
\[
R(U,V) =
 \sum_{r=1}^m\sigma_r\ell_r
+ 2\sigma_m\tau.
\]
By the triangle inequality,
\[
|R(U,V)|
 \le
 \sum_{r=1}^m\ell_r+2\tau
= |R(U^\sharp,V^\sharp)|.
\]
Since all \(\ell_r\) and \(\tau\) are positive, equality holds if and only if
\[
 \sigma_1=\sigma_2=\cdots=\sigma_m.
\]
The two possible common signs simply interchange
\(U^\sharp\) and \(V^\sharp\).
Hence the unique unordered pair minimizing $Q$ and then maximizing $|R|$ is
\[
 \left\{
 I_s\cup\{d_{2r-1}:1\le r\le m\}\cup\{d_{2m}\},
 \;
 I_s\cup\{d_{2r}:1\le r\le m-1\}
\right\}.
\]
This is the even-sorted pair, so $q_{U^\sharp}q_{V^\sharp}$ is the unique quadratic minimizer in Case~2.

\medskip

Therefore the assertion follows in both cases.
\end{proof}

\end{document}
